\section{Regularized derivatives and odd spectral values}
\label{ix-add:bendersky}

The continuation and functional equation proved in
Proposition~\ref{rz-num-03-2} relate two operations on the same
functional $Z$: differentiation at negative integers and convergent
evaluation at positive odd integers. The regional coordinates and
enumeration operator retain their preceding genealogy. The symbol
$\pi$ denotes the HMT coordinate already used in the theta--Mellin
construction.

Let $H_k=\sum_{j=1}^k1/j$, with $H_0=0$, and let $\mathsf B_j$
be the Bernoulli numbers constructed by the recurrence in
Section~\ref{lectura:manuscrito-04}, with $\mathsf B_1=-1/2$.
The notation $\mathsf B_j$ distinguishes these coefficients from the
Meissel--Mertens constant. Define the Bendersky--Glaisher--Kinkelin
hierarchy by
\[
 A_k=\exp\!\left(\frac{H_k\mathsf B_{k+1}}{k+1}-Z'(-k)\right),
 \qquad k\in\mathbb N_0.
\]
The points $-k$ are regular. The Bernoulli term fixes the normalization
of each level before evaluation.

\subsection{Covariance of regularized derivatives}
Here the functional is $Z_\lambda=\lambda^{-s}Z$, with residue
$\lambda^{-1}$, and we define
\[
 D_k(\lambda)=\exp[-Z_\lambda'(-k)].
\]
This object retains the normalization of the scaled operator, distinct
from the unit-residue normalization used for $\widehat Z_\lambda$
in Section~\ref{ix-add:stieltjes}.

\HMTresultado{Proposition}{Transport of regularized derivatives}{ix-add:regularized-scaling}
For $\lambda>0$ and $k\in\mathbb N_0$,
\[
 Z_\lambda'(-k)=\lambda^k
       \bigl(Z'(-k)-(\log\lambda)Z(-k)\bigr),
\]
\[
 D_k(\lambda)=D_k(1)^{\lambda^k}
               \lambda^{\lambda^k Z(-k)},\qquad
 D_k(1)=A_k\exp\!\left(-\frac{H_k\mathsf B_{k+1}}{k+1}\right).
\]
\textbf{Proof.} The product rule applied to $\lambda^{-s}Z(s)$ gives
the first identity. Its negative exponential gives the second, and
the definition of $A_k$ gives the third. \hfill$\square$

At the initial levels,
\[
\begin{aligned}
 A_0&=e^{-Z'(0)},\\
 \log A_1&=\frac1{12}-Z'(-1),\qquad \log A_2=-Z'(-2),\\
 \log A_3&=-\frac{11}{720}-Z'(-3),\qquad \log A_4=-Z'(-4).
\end{aligned}
\]
For example, $D_1(1)=A_1e^{-1/12}$. Retaining the normalization
prevents two different readouts from being identified merely because
they involve the same derivative.

\Needspace{12\baselineskip}
\subsection{Recovery of odd values}
\HMTresultado{Theorem}{Correspondence between negative even and positive odd levels}{ix-add:odd-values}
For every integer $n\geq1$,
\[
 Z'(-2n)=(-1)^n\frac{(2n)!}{2(2\pi)^{2n}}Z(2n+1),
\]
\[
 \frac{\log D_{2n}(\lambda)}{\lambda^{2n}}
 =\log A_{2n}
 =(-1)^{n+1}\frac{(2n)!}{2(2\pi)^{2n}}Z(2n+1).
\]
\textbf{Proof.} The functional equation of the completed form implies
the meromorphic identity
\[
 \pi^{-s/2}\Gamma(s/2)Z(s)
 =\pi^{-(1-s)/2}\Gamma((1-s)/2)Z(1-s).
\]
At $s=-2n$, the right-hand side is regular and the gamma factor on
the left has a simple pole. Hence $Z(-2n)=0$. The gamma recurrence
and its residue at zero give
\[
 \Gamma(s/2)=\frac{2(-1)^n}{n!(s+2n)}+O(1).
\]
Moreover,
\[
 \Gamma\!\left(n+\frac12\right)
 =\frac{(2n)!\sqrt\pi}{4^nn!}.
\]
Comparing constant terms in the meromorphic identity yields
\[
 \pi^n\frac{2(-1)^n}{n!}Z'(-2n)
 =\pi^{-n}\frac{(2n)!}{4^nn!}Z(2n+1),
\]
which gives the first formula. The series $t/(e^t-1)+t/2$ is even;
therefore $\mathsf B_{2n+1}=0$ for $n\geq1$. The definition of
$A_{2n}$ and the preceding proposition complete the proof.
\hfill$\square$

In particular, the first positive even level satisfies
\[
 Z(3)=4\pi^2\log A_2.
\]
The Ap\'ery constant evaluated with a remainder in
Section~\ref{lectura:manuscrito-04} is thus related to a regularized
derivative of the same functional. The intervals proved there are
transported by the monotone exponential once the bounds for the
regional coordinate $\pi$ are retained.

Level zero has a different evaluation. The functional equation,
together with gamma reflection and duplication, is equivalent to
\[
 Z(s)=2^s\pi^{s-1}\sin(\pi s/2)\Gamma(1-s)Z(1-s).
\]
Euler's formula
\[
 \Gamma(1+z)=\lim_{m\to\infty}
 m^z\prod_{j=1}^{m}(1+z/j)^{-1}
\]
gives $\log\Gamma(1+z)=-\gamma_0z+O(z^2)$: the linear coefficient
is the limit of $\log m-H_m$, and the quadratic remainder is controlled
by the convergent series of $j^{-2}$. Using
$Z(1-s)=-1/s+\gamma_0+O(s)$ and
$\Gamma(1-s)=1+\gamma_0s+O(s^2)$, the degree-one terms involving
$\gamma_0$ cancel, giving
\[
 Z(s)=-\frac12-\frac12\log(2\pi)s+O(s^2).
\]
Consequently,
\[
 A_0=\sqrt{2\pi},\qquad D_0(\lambda)=\sqrt{\frac{2\pi}{\lambda}}.
\]
The condition $n\geq1$ in the theorem preserves this distinction:
level zero is not obtained by substituting $Z(1)$, which has a pole,
into the odd-value formula. The gamma identities express subsequent
analytic normalization without participating in seed selection.
